\section{Introduction}
There are multiple methods by which to incorporate additional respiratory phase information into image reconstruction and they differ in how they approach handling the data. The iMoCo and XD-GRASP methods are the two primary methods of interest in this proposal and both are novel in their application to 4D flow of the portal venous system28,29. However, prior studies in the portal venous system using 4D CT and 2D PC MRI have shown that there can be respiration-induced differences in portal venous flow30,31. 

\section{Materials \& Methods}
Data acquired for Aim 2.1 will be used for this subaim as well. The PCVIPR sequence can be reconstructed in different respiratory states similar to what was described in Aim 1.2. I will reconstruct a subset of the acquired data (n=10, venc=30 cm/s) with a 50\% inspiration window and reanalyze the datasets as described in Aim 2.1. If the previously described 4D Laplacian method is successful in unwrapping aliased phases, I will perform the analysis with the phase unwrapped data, otherwise, I will use the data unwrapped with the dual-venc method. I will then compare flow differences between inspiration and expiration to establish maximum deviations across the dataset and compare conservation of mass as a consistency measurement.

It is expected that the differences in peak flow between respiratory phases will be relatively small, so I intend to implement iMoCo (iterative motion correction) to treat extra respiratory phases as motion-corrupted data and incorporate the information back into a single respiratory phase (expiration)28. The respiratory bellows signal will be used to perform a low-res, respiratory-resolved reconstruction, sorting data into 8\-10 respiratory bins to minimize residual motion within each phase. All respiratory bins will then be registered to the end-expiration reference phase via deformable image registration44. The motion fields for each phase are then included as a linear operator in the compressed sensing reconstruction. 

To test the effectiveness of the iMoCo approach, acquired data will be reconstructed with both the standard 50\% expiration window and the iMoCo method. All sets of images will be reanalyzed in GTFlow reusing the same plane placements done in Aim 2.1. Correlation of vessel velocities and conservation of mass will be compared between each reconstruction as measures of precision and repeatability. Image quality will be assessed by a reader study with an experienced radiologist scoring the quality of the images according to the Likert scale. To test the feasible reduction in scan time, the iMoCo acquisitions will also be downsampled to use 50\% and 25\% of the acquired data (a theoretical scan time reduction to 6 and 3 min) and then compared with the standard 50\% expiration window.


\section{Results}

\section{Discussion}

\section{Conclusion}